\section{Introduction}
\label{sec:introduction}

The evolution toward 6G networks is driven by the need to support increasingly
complex and dynamic applications, including Digital Twins (DT), autonomous
systems, and intelligent environments. Among these, smart-warehouse systems
represent a particularly challenging use case, where mobile robots, sensing
devices, and control systems must operate reliably under stringent latency,
reliability, and throughput requirements. In such environments, the wireless
network is no longer a passive communication medium: it becomes an active
component of the system that must adapt continuously to changes in user
position, mobility patterns, and environmental conditions. This motivates the
concept of \emph{situation awareness}, in which the communication system
leverages contextual information to optimise its behaviour proactively rather
than reactively~\cite{ref:positioning_6g}.

Situation awareness, in the context of wireless access networks, refers to the
ability of the network---particularly the antenna system and its associated
control logic---to perceive and exploit contextual information such as user
equipment (UE) location, mobility trajectories, surrounding geometry, and
potential obstacles. Instead of relying on periodic beam training or
threshold-based handovers, a situation-aware system anticipates future channel
conditions and adapts accordingly. In a smart-warehouse setting, this
includes predicting robot movement along planned paths, identifying potential
blockages from racks or moving objects, and adjusting beamforming strategies
to maintain reliable links. Such proactive adaptation has been shown to
improve key performance indicators (KPIs) including latency, throughput,
reliability, and energy efficiency, particularly in higher-frequency bands
where propagation is highly sensitive to environmental
dynamics~\cite{ref:multimodal_bf,ref:predictive_handover}.

Enabling situation awareness, however, requires the network to have access to
accurate and timely environmental information. Conventional radio measurements
such as RSRP or SINR provide limited insight into the physical environment and
often arrive too late to support proactive adaptation. To overcome this
limitation, the antenna system must be endowed with \emph{sensing
capabilities}, allowing it to observe the surrounding environment and infer
the position and motion of objects directly. This convergence of sensing and
communication functionalities is a fundamental requirement for next-generation
networks and forms the basis of the Integrated Sensing and Communication
(ISAC) paradigm. ISAC integrates radar-like sensing into the communication
infrastructure so that the same hardware, waveform, and spectrum are jointly
used for data transmission and environmental sensing. In a smart warehouse,
ISAC enables a base station to detect and track mobile robots and other
dynamic objects from reflected signals, transforming the communication system
into a distributed sensing platform without dedicated sensing infrastructure.

A key challenge in studying ISAC systems lies in faithfully modelling the
interaction between electromagnetic waves and complex environments.
Conventional system-level simulators rely on stochastic channel models that do
not capture site-specific reflections, diffractions, and blockages. To address
this, ray-tracing techniques are employed as a high-fidelity modelling
approach: by leveraging detailed scene geometry and material properties, ray
tracing provides per-path delays, angles of arrival/departure, and complex
gains, which are essential both for accurate communication modelling and for
the radar-like reasoning that underpins
sensing~\cite{ref:sionna_rt,ref:ns3sionna}.

Building on this foundation, this paper presents an ISAC framework that pairs
NVIDIA Sionna RT with ns-3/5G~LENA into a single co-simulation pipeline, and
evaluates it both in controlled scenes and in an end-to-end smart-warehouse
Digital Twin. The contributions of the paper are as follows:
\begin{itemize}
  \item It describes the integration of ns-3/5G~LENA with Sionna RT through a
        ZeroMQ-based propagation/sensing bridge, with a clean role boundary
        between the two simulators (Section~\ref{sec:methodology}).
  \item It presents a four-module ISAC architecture comprising a
        radar-inspired sensing module (Equivalent Reflection Point (ERP)
        computation, $Z$-gating, Moving-Target Indication (MTI), and DBSCAN
        clustering), a Kalman--Hungarian multi-target tracking module, a
        sensing-informed multi-lobe adaptive beamforming module, and a
        communication module that exports per-path channel state back to
        ns-3 (Section~\ref{sec:methodology}).
  \item It validates the framework through a two-stage evaluation: a
        controlled scene-level study across free-space, warehouse, and
        urban-micro street-canyon scenes with six ablation variants, and an
        application-oriented warehouse Digital Twin study under three gNB
        array sizes (Section~\ref{sec:results}).
  \item It identifies the operating regimes in which the proposed framework
        is beneficial, including the $4{\times}4$ gNB configuration as the
        balanced operating point in the smart-warehouse scenario, and
        discusses the remaining limitations and future research directions
        (Sections~\ref{sec:results}--\ref{sec:conclusion}).
\end{itemize}

The remainder of this paper is organised as follows.
Section~\ref{sec:methodology} describes the simulation flow, ISAC
architecture, and the algorithms used in each module.
Section~\ref{sec:results} presents the controlled-validation and
warehouse-use-case results together with a discussion of their implications.
Section~\ref{sec:conclusion} concludes the paper and outlines future work.
